\documentclass[12pt,a4paper]{article}
\usepackage[a4paper,left=2.5cm,right=2cm,top=2cm,bottom=2cm]{geometry}
\usepackage{fontspec}
\setmainfont{DejaVu Serif}
\setmonofont{DejaVu Sans Mono}
\usepackage{setspace}
\onehalfspacing
\usepackage{enumitem,booktabs,tabularx,array,xcolor,listings,hyperref,titlesec}
\hypersetup{hidelinks}
\titleformat{\section}{\normalfont\Large\bfseries}{\thesection.}{0.5em}{}
\titleformat{\subsection}{\normalfont\large\bfseries}{\thesubsection.}{0.5em}{}
\setlength{\parindent}{1.25cm}
\setlength{\parskip}{0pt}
\definecolor{codebg}{RGB}{246,246,246}
\definecolor{codeframe}{RGB}{215,215,215}
\lstset{basicstyle=\ttfamily\small,backgroundcolor=\color{codebg},frame=single,rulecolor=\color{codeframe},breaklines=true,columns=fullflexible,keepspaces=true,showstringspaces=false,aboveskip=.8em,belowskip=.8em}
\begin{document}
\begin{titlepage}\begin{center}
{\large Ministerul Educației și Cercetării al Republicii Moldova}\\[.3cm]
{\large Universitatea Tehnică a Moldovei}\\[.3cm]
{\large Facultatea Calculatoare, Informatică și Microelectronică}\\[.3cm]
{\large Program: Software Engineering}\\[2.3cm]
{\Large\textbf{Operating Systems}}\\[.8cm]
{\LARGE\textbf{Laboratory Work No. 10}}\\[.5cm]
{\Large\textbf{Add Your Own System Call: trace}}\\[.35cm]
{\large Track B: build your own small OS (xv6)}
\vfill
\begin{flushright}\begin{tabular}{ll}
\textbf{Elaborated by:} & \\
st. gr. FAF-241 & \\
\textbf{Ermicev Sofia} & \\
\\
\textbf{Verified by:} & \\
\textbf{Patricia Reițman} &
\end{tabular}\end{flushright}
\vfill {\large Chișinău -- 2026}
\end{center}\end{titlepage}
\tableofcontents\newpage
\section{Introduction}
The aim of this laboratory work is to extend the xv6 RISC-V operating system with a new system call named \texttt{trace}. Unlike Lab 9, where \texttt{strace} observed calls from outside a Linux process, this work modifies the xv6 user interface and kernel to report selected system calls. The experiment uses system call number 23 and a per-process bit mask. Work was performed in the existing \texttt{/home/sofia/xv6-riscv} tree on Ubuntu under WSL2.

\section{Part 1. The Number and the User Side}
\subsection{Registration}
The value \texttt{SYS\_trace = 23} was added after \texttt{SYS\_sync = 22} in \texttt{kernel/syscall.h}. The user-level declaration was added to \texttt{user/user.h}, and \texttt{user/usys.pl} was extended so that the generated stub can issue the RISC-V \texttt{ecall} instruction.
\begin{lstlisting}
// kernel/syscall.h
#define SYS_sync   22
#define SYS_trace  23

// user/user.h
int sync(void);
int trace(int);

# user/usys.pl
entry("sync");
entry("trace");
\end{lstlisting}
\subsection{User program}
A new \texttt{user/trace.c} accepts a mask and a command, calls \texttt{trace(mask)}, and replaces itself with the requested program using \texttt{exec()}. The program was added to the \texttt{UPROGS} list in the Makefile.
\begin{lstlisting}[language=C]
#include "kernel/types.h"
#include "kernel/stat.h"
#include "user/user.h"
int main(int argc, char *argv[])
{
  if(argc < 3){
    fprintf(2, "usage: trace mask command [args...]\n");
    exit(1);
  }
  if(trace(atoi(argv[1])) < 0){
    fprintf(2, "trace: failed\n");
    exit(1);
  }
  exec(argv[2], argv + 2);
  fprintf(2, "trace: exec %s failed\n", argv[2]);
  exit(1);
}
\end{lstlisting}
\subsection{Check 1: Unimplemented kernel call}
Before implementing the handler, xv6 booted and produced the following output.
\begin{lstlisting}
$ trace 32 grep hello README
3 trace: unknown sys call 23
trace: failed
\end{lstlisting}
This was the intended intermediate result: the user program issued system call 23, but the kernel had no registered handler yet.

\section{Part 2. The Kernel Side}
\subsection{Process mask and handler}
The \texttt{trace\_mask} field was added after the process name in \texttt{struct proc} (\texttt{kernel/proc.h}). The handler in \texttt{kernel/sysproc.c} reads the integer argument and stores it in the current process.
\begin{lstlisting}[language=C]
// kernel/proc.h, inside struct proc
char name[16];
int trace_mask;

// kernel/sysproc.c
uint64
sys_trace(void)
{
    int mask;
    argint(0, &mask);
    myproc()->trace_mask = mask;
    return 0;
}
\end{lstlisting}
\subsection{Dispatch, output, and inheritance}
The handler was declared and added to the dispatcher in \texttt{kernel/syscall.c}. A name table maps call numbers to readable names. After a selected system call finishes, the kernel prints its PID, name, and return value. This xv6 version uses \texttt{printk()} rather than \texttt{printf()} for kernel output.
\begin{lstlisting}[language=C]
extern uint64 sys_trace(void);
// In syscalls[]:
[SYS_trace] = sys_trace,

// Immediately after p->trapframe->a0 = syscalls[num]();
if ((p->trace_mask >> num) & 1)
    printk("%d: syscall %s -> %ld\n",
           p->pid, syscallnames[num], p->trapframe->a0);

// kernel/proc.c, in kfork(), after copying the name:
np->trace_mask = p->trace_mask;
\end{lstlisting}
The inheritance line ensures that children created with \texttt{fork()} keep the parent's tracing mask.
\subsection{Check 2: Tracing read}
After recompiling and booting xv6, the specified command produced:
\begin{lstlisting}
$ trace 32 grep hello README
3: syscall read -> 1023
3: syscall read -> 965
3: syscall read -> 453
3: syscall read -> 0
\end{lstlisting}
The successful reads total \(1023+965+453=2441\) bytes, equal to the README size stated in the laboratory sheet. The final return value of zero indicates end of file.

\section{Part 3. Playing with the Mask}
The mask uses bit \(N\) for system call number \(N\). The experiments below were run in xv6; PIDs are taken from the terminal output.
\subsection{All calls}
The mask \texttt{2147483647} enables bits 0 through 30. The observed output was:
\begin{lstlisting}
$ trace 2147483647 grep hello README
4: syscall trace -> 0
4: syscall exec -> 3
4: syscall open -> 3
4: syscall read -> 1023
4: syscall read -> 965
4: syscall read -> 453
4: syscall read -> 0
4: syscall close -> 0
\end{lstlisting}
\subsection{Only fork}
With mask 2, the \texttt{forkforkfork} test displayed fork calls from several PIDs and finished successfully. Selected lines from the actual output:
\begin{lstlisting}
$ trace 2 usertests forkforkfork
usertests starting
test forkforkfork: 5: syscall fork -> 6
6: syscall fork -> 7
7: syscall fork -> 8
7: syscall fork -> 9
...
8: syscall fork -> 66
8: syscall fork -> -1
7: syscall fork -> -1
9: syscall fork -> -1
10: syscall fork -> -1
OK
ALL TESTS PASSED
\end{lstlisting}
The negative return values indicate unsuccessful attempts once resources were exhausted in this test.
\subsection{Only write}
\begin{lstlisting}
$ trace 65536 echo hi
hi67: syscall write -> 2

67: syscall write -> 1
\end{lstlisting}
The \texttt{echo} program made one write of two bytes for \texttt{hi} and another of one byte for the newline. The kernel logged the first call before the second write produced the newline.
\subsection{Only open and close}
\begin{lstlisting}
$ trace 2129920 grep hello README
68: syscall open -> 3
68: syscall close -> 0
\end{lstlisting}
The exact mask is \(2^{15}+2^{21}=32768+2097152=2129920\).

\subsection{Observations}
\begin{enumerate}[leftmargin=*]
\item \textbf{Why does mask 32 show only read?} The mask 32 equals \(2^5\), so only bit 5 is set. Since \texttt{SYS\_read} has number 5, \texttt{open} and \texttt{close} are not selected.
\item \textbf{Why do multiple PIDs print fork lines?} The line \texttt{np->trace\_mask = p->trace\_mask;} in \texttt{kfork()} copies the mask into every child process. This was observable in the PIDs printed by the test.
\item \textbf{Why is tracing printed between hi and the newline?} \texttt{echo} performs two \texttt{write()} calls: one for the two characters and one for the newline. The trace message is printed immediately after each call returns.
\item \textbf{Which mask selects only open and close?} \(2^{15}+2^{21}=2129920\), as verified by the last experiment.
\end{enumerate}

\section{Conclusion}
I extended xv6 with a user-facing system call and a kernel-side handler, connected the call number to the dispatch table, and used a per-process bitmask to control printed traces. The initial unknown-system-call result confirmed that the user-to-kernel interface was functioning; the subsequent read trace confirmed the full implementation. The fork test demonstrated why trace settings must be inherited by child processes. Overall, this exercise connected the system-call observations from Linux Lab 9 to the implementation of the same mechanism inside xv6.
\end{document}
